\chapter{Заключение}

В рамках данной работы был проведён подробный анализ ранее предложенного стохастического алгоритма робастного оценивания многомерного параметра положения. Основное внимание было уделено исследованию теоретических свойств алгоритма, а также разработке модификаций, направленных на повышение точности и устойчивости оценки.

Была установлена важная теоретическая связь предложенного алгоритма с проекционной медианой и методом yamm. Показано, что получаемая алгоритмом оценка эквивалентна оценке по методу yamm, что позволяет перенести на неё известные теоретические свойства проекционной медианы, включая эквивариантность относительно вращений.

Проведено исследование эквивариантности алгоритма. Показано, что для ортогональных преобразований отклонение от эквивариантности оказывается близким к нулю, тогда как для более общих аффинных преобразований свойство аффинной эквивариантности не выполняется. При этом характер поведения алгоритма оказался сопоставим с геометрической медианой, которая являлась основным конкурирующим методом в рамках работы.

Для повышения устойчивости и точности оценки были предложены и исследованы несколько модификаций алгоритма. Рассматривались условия на минимальное число итераций, многократное выполнение критерия останова, усреднение последних итераций, экспоненциальное сглаживание, а также адаптивный выбор шага на основе метода AdaGrad. Наиболее устойчивые результаты были получены для модификаций, использующих экспоненциальное сглаживание и усреднение последних итераций. Проведённые численные эксперименты показали, что данные подходы позволяют существенно уменьшить разброс оценки и снизить влияние случайности, возникающей из-за стохастической природы алгоритма.

Кроме того, были исследованы специальные примеры, в которых базовый алгоритм не демонстрирует устойчивой сходимости. Рассматривались выборки в форме букв V и C, а также выборки, состоящие из вершин вложенных треугольников. Показано, что предложенные модификации позволяют добиться приближённой сходимости даже в подобных сложных случаях: траектории алгоритма концентрируются вблизи истинного параметра положения и образуют компактные области вместо хаотичного поведения базовой версии алгоритма.

Отдельное внимание было уделено исследованию поведения алгоритма на выборках различной природы. Проведены эксперименты для непрерывных и дискретных распределений, а также для выборок, заданных в выпуклых, невыпуклых и несвязных областях. Результаты показали, что предложенный алгоритм остаётся применимым и в случаях, не связанных с эллиптически контурированными распределениями, хотя теоретическое объяснение данного поведения требует дальнейшего исследования.

В работе также была рассмотрена задача робастного оценивания матрицы разброса. Предложен метод, основанный на использовании случайных направлений и робастных оценок масштаба вдоль проекций выборки. Было показано, что данный подход связан с формулой Гнанадесикана--Кеттеринга, однако позволяет по-новому интерпретировать задачу восстановления матрицы разброса через случайные проекции. Кроме того, рассмотрен альтернативный вычислительно устойчивый подход, позволяющий избежать проблем, связанных с плохой обусловленностью возникающих систем линейных уравнений.

Таким образом, в работе удалось не только исследовать свойства исходного стохастического алгоритма, но и предложить ряд его улучшений, существенно повышающих устойчивость и качество оценивания. Полученные результаты показывают перспективность подходов, основанных на случайных проекциях и робастных одномерных статистиках, для задач многомерного робастного анализа данных.

В качестве направлений дальнейших исследований представляют интерес теоретическое изучение сходимости алгоритма для неэллиптических распределений, исследование аффинно-эквивариантных модификаций метода, а также дальнейшее развитие подходов к робастному оцениванию матрицы разброса.



% В рамках работы был проведен более подробный анализ предложенного ранее робастного алгоритма для оценки многомерного параметра положения. Особое внимание уделено его модификациям, направленным на повышение устойчивости и точности. Одной из ключевых доработок стало добавление экспоненциального сглаживания с усреднением последних итераций, что позволило значительно снизить ошибку оценки и повысить стабильность алгоритма.

% % Важным результатом стало эмпирическое исследование эквивариантности алгоритма. Для ортогональных преобразований отклонение от эквивариантности оказалось близким к нулю, а для диагональных матриц и матриц сдвига медианная ошибка не превышала $0,1$, что сопоставимо с точностью геометрической медианы. Однако эквивариантности для произвольной матрицы преобразования не наблюдается.

% Была получена важная теоретическая связь оценки, полученной предложенным алгоритмом, с оценкой проекционной медианы и оценки по методу yamm. Показана эквивалентность полученной оценки с оценкой по методу yamm, из чего сделан вывод о совпадении теоретических свойств оценки.

% Помимо этого, были рассмотрены случаи, когда алгоритм не сходится, которые демонстрируют, что некоторые специально подобранные примеры являются сложными задачами для предложенного алгоритма. Однако для таких проблемных случаев (выборки в форме V, C, вершин треугольников) предложенная модификация с экспоненциальным сглаживанием и усреднением обеспечивает приближенную сходимость, когда оценки образуют компактное облако вблизи истинного значения.

% Кроме того, была проведена эмпирическая оценка точности алгоритма для различных типов областей и распределений, показавшая его применимость для распределений, не являющихся эллиптически контурированными. Однако теоретически пока сложно доказать, почему предложенный алгоритм в данных случаях сходится к истинному значению параметра положения.

% Дополнительно была разработана робастная оценка матрицы разброса, что стало важным шагом для полного описания многомерных данных. Было проведено улучшение методов оценки матриц разброса. Показано, что предложенная оценка имеет связь с ранее опубликованной оценкой, которая, однако, не устойчива относительно вращений данных. Приведен новый метод, который устраняет недостатки ранее предложенных оценок.

% Также рассмотрены методы улучшения сходимости последовательности при оценивании параметра положения. Разность между двумя соседними оценками в последовательности была интепретирована, как направление поиска медианы, после чего, помимо экспоненциального сглаживания, удалось рассмотреть модификацию для улучшения сходимости AdaGrad, позаимствованную из метода стохастического градиента.

% На основе проделанной работы определены перспективные направления для дальнейших исследований.